\origpage{56}{ROBERTSON AND SEYMOUR}
\noindent by each of $S_{j-\nu+1}, S_{j-\nu+2},..., S_j$, by (4.1). Then we have a contradiction
from (2.7).

For $\nu\le j\le n$, choose $X_j$ as in (4.2), minimal.

For $1\le j\le n$, let $Z_j$ denote the set of all vertices of $G$ which can be joined
to a vertex of $S_j$ by a path of $G$ which avoids $(S_1\cup\cdots\cup S_n)-S_j$. Then
clearly $Z_j\cap(S_1\cup\cdots\cup S_n)=S_j$ and $Z_1\cup\cdots\cup Z_n=V(G)$. By (4.1),
$Z_j\cap Z_{j'}\ne\varnothing$ only if $|j-j'|\le 1$.
\stmt
(4.3). \textit{For $\nu\le j\le n$, $X_j\subseteq S_{j-\nu+1}\cup Z_{j-\nu+2}\cup\cdots\cup Z_{j-1}\cup S_j$.}
\stmt
For if $v\in X_j$, then by the minimality of $X_j$ there is a path $P$ from
$\bigcup_{1\le j'\le j-\nu+1}S_{j'}$ to $\bigcup_{j\le j'\le n}S_{j'}$ which avoids $X_j-\{v\}$ and uses $v$, and which
has no interior vertices in
\begin{equation*}
\bigcup_{1\le j'\le j-\nu+1}S_{j'}\cup\bigcup_{j\le j'\le n}S_{j'}.
\end{equation*}
Since $S_{j-\nu+2}\cup\cdots\cup S_{j-1}$ separates $\bigcup_{1\le j'\le j-\nu+1}S_{j'}$ and $\bigcup_{j\le j'\le n}S_{j'}$ (by
(4.1)), $P$ has an interior vertex $u$ in $S_{j-\nu+2}\cup\cdots\cup S_{j-1}$. Choose $u$ such
that the subpath $P'$ of $P$ from $v$ to $u$ has minimum length. Then no vertices
of $P'$ except $v$, $u$ are in $S_1\cup\cdots\cup S_n$. If $v\notin S_1\cup\cdots\cup S_n$, then
$v\in Z_{j-\nu+2}\cup\cdots\cup Z_{j-1}$, because of $P'$. If $v\in S_1\cup\cdots\cup S_n$, then
\begin{equation*}
v\in S_{j-\nu+1}\cup S_{j-\nu+2}\cup\cdots\cup S_{j-1}\cup S_j
\end{equation*}
by (4.1), since no interior vertices of $P'$ are in $S_{j-\nu+1}$ or $S_j$. Thus in either
case,
\begin{equation*}
v\in S_{j-\nu+1}\cup Z_{j-\nu+2}\cup\cdots\cup Z_{j-1}\cup S_j
\end{equation*}
as claimed.
\stmt
(4.4). \textit{For $\nu\le i\le n$, $1\le j\le n$, if $s\in S_j$, then}

\hspace*{12pt}(i)\enspace\textit{$X_i$ separates $s$ and $\rho(G)$ if $j\ge \delm{1}\fix{i}$, and}\fixnote{fix:4.4}{As printed, (i) and (ii) contradict each other for $i=n$ and $s=\rho(G)\in S_1$, since $1\le n-\nu$, and $\rho(G)\notin X_n$ by (4.3). The proof of (i) works for $j\ge i$: $X_i$ separates $S_1\cup\dots\cup S_{i-\nu+1}$, which contains $\rho(G)$, from $S_i\cup\dots\cup S_n$, which contains $s$. The paper uses (i) only for such $j$, in (4.5) and (4.6).}

\hspace*{12pt}(ii)\enspace\textit{$X_i$ does not separate $s$ and $\rho(G)$ if $j\le i-\nu$.}
\stmt
Statement (i) follows from the definition of $X_i$, since $\rho(G)\in S_1$. To show
(ii), we observe that from (4.3), the path of $H$ from $s$ to $\rho(G)$ does not meet
$X_i$ if $j\le i-\nu$.

Let $Y_i=X_{\nu i}$ for $1\le i\le T$. From (4.3), $Y_1,..., Y_T$ are disjoint.
\stmt
(4.5). \textit{For $1\le i, i'\le T$ and for every $y\in Y_{i'}$, $Y_i$ separates $y$ and $\rho(G)$ if
and only if $i\le i'$.}
